% Options for packages loaded elsewhere
\PassOptionsToPackage{unicode}{hyperref}
\PassOptionsToPackage{hyphens}{url}
\documentclass[
]{article}
\usepackage{xcolor}
\usepackage{amsmath,amssymb}
\setcounter{secnumdepth}{-\maxdimen} % remove section numbering
\usepackage{iftex}
\ifPDFTeX
  \usepackage[T1]{fontenc}
  \usepackage[utf8]{inputenc}
  \usepackage{textcomp} % provide euro and other symbols
\else % if luatex or xetex
  \usepackage{unicode-math} % this also loads fontspec
  \defaultfontfeatures{Scale=MatchLowercase}
  \defaultfontfeatures[\rmfamily]{Ligatures=TeX,Scale=1}
\fi
\usepackage{lmodern}
\ifPDFTeX\else
  % xetex/luatex font selection
\fi
% Use upquote if available, for straight quotes in verbatim environments
\IfFileExists{upquote.sty}{\usepackage{upquote}}{}
\IfFileExists{microtype.sty}{% use microtype if available
  \usepackage[]{microtype}
  \UseMicrotypeSet[protrusion]{basicmath} % disable protrusion for tt fonts
}{}
\makeatletter
\@ifundefined{KOMAClassName}{% if non-KOMA class
  \IfFileExists{parskip.sty}{%
    \usepackage{parskip}
  }{% else
    \setlength{\parindent}{0pt}
    \setlength{\parskip}{6pt plus 2pt minus 1pt}}
}{% if KOMA class
  \KOMAoptions{parskip=half}}
\makeatother
\usepackage{color}
\usepackage{fancyvrb}
\newcommand{\VerbBar}{|}
\newcommand{\VERB}{\Verb[commandchars=\\\{\}]}
\DefineVerbatimEnvironment{Highlighting}{Verbatim}{commandchars=\\\{\}}
% Add ',fontsize=\small' for more characters per line
\newenvironment{Shaded}{}{}
\newcommand{\AlertTok}[1]{\textcolor[rgb]{1.00,0.00,0.00}{\textbf{#1}}}
\newcommand{\AnnotationTok}[1]{\textcolor[rgb]{0.38,0.63,0.69}{\textbf{\textit{#1}}}}
\newcommand{\AttributeTok}[1]{\textcolor[rgb]{0.49,0.56,0.16}{#1}}
\newcommand{\BaseNTok}[1]{\textcolor[rgb]{0.25,0.63,0.44}{#1}}
\newcommand{\BuiltInTok}[1]{\textcolor[rgb]{0.00,0.50,0.00}{#1}}
\newcommand{\CharTok}[1]{\textcolor[rgb]{0.25,0.44,0.63}{#1}}
\newcommand{\CommentTok}[1]{\textcolor[rgb]{0.38,0.63,0.69}{\textit{#1}}}
\newcommand{\CommentVarTok}[1]{\textcolor[rgb]{0.38,0.63,0.69}{\textbf{\textit{#1}}}}
\newcommand{\ConstantTok}[1]{\textcolor[rgb]{0.53,0.00,0.00}{#1}}
\newcommand{\ControlFlowTok}[1]{\textcolor[rgb]{0.00,0.44,0.13}{\textbf{#1}}}
\newcommand{\DataTypeTok}[1]{\textcolor[rgb]{0.56,0.13,0.00}{#1}}
\newcommand{\DecValTok}[1]{\textcolor[rgb]{0.25,0.63,0.44}{#1}}
\newcommand{\DocumentationTok}[1]{\textcolor[rgb]{0.73,0.13,0.13}{\textit{#1}}}
\newcommand{\ErrorTok}[1]{\textcolor[rgb]{1.00,0.00,0.00}{\textbf{#1}}}
\newcommand{\ExtensionTok}[1]{#1}
\newcommand{\FloatTok}[1]{\textcolor[rgb]{0.25,0.63,0.44}{#1}}
\newcommand{\FunctionTok}[1]{\textcolor[rgb]{0.02,0.16,0.49}{#1}}
\newcommand{\ImportTok}[1]{\textcolor[rgb]{0.00,0.50,0.00}{\textbf{#1}}}
\newcommand{\InformationTok}[1]{\textcolor[rgb]{0.38,0.63,0.69}{\textbf{\textit{#1}}}}
\newcommand{\KeywordTok}[1]{\textcolor[rgb]{0.00,0.44,0.13}{\textbf{#1}}}
\newcommand{\NormalTok}[1]{#1}
\newcommand{\OperatorTok}[1]{\textcolor[rgb]{0.40,0.40,0.40}{#1}}
\newcommand{\OtherTok}[1]{\textcolor[rgb]{0.00,0.44,0.13}{#1}}
\newcommand{\PreprocessorTok}[1]{\textcolor[rgb]{0.74,0.48,0.00}{#1}}
\newcommand{\RegionMarkerTok}[1]{#1}
\newcommand{\SpecialCharTok}[1]{\textcolor[rgb]{0.25,0.44,0.63}{#1}}
\newcommand{\SpecialStringTok}[1]{\textcolor[rgb]{0.73,0.40,0.53}{#1}}
\newcommand{\StringTok}[1]{\textcolor[rgb]{0.25,0.44,0.63}{#1}}
\newcommand{\VariableTok}[1]{\textcolor[rgb]{0.10,0.09,0.49}{#1}}
\newcommand{\VerbatimStringTok}[1]{\textcolor[rgb]{0.25,0.44,0.63}{#1}}
\newcommand{\WarningTok}[1]{\textcolor[rgb]{0.38,0.63,0.69}{\textbf{\textit{#1}}}}
\usepackage{longtable,booktabs,array}
\usepackage{caption}
\captionsetup[table]{skip=6pt}
\newcounter{none} % for unnumbered tables
\usepackage{calc} % for calculating minipage widths
% Correct order of tables after \paragraph or \subparagraph
\usepackage{etoolbox}
\makeatletter
\patchcmd\longtable{\par}{\if@noskipsec\mbox{}\fi\par}{}{}
\makeatother
\usepackage{graphicx}
\makeatletter
\newsavebox\pandoc@box
\newcommand*\pandocbounded[1]{% scales image to fit in text height/width
  \sbox\pandoc@box{#1}%
  \Gscale@div\@tempa{\textheight}{\dimexpr\ht\pandoc@box+\dp\pandoc@box\relax}%
  \Gscale@div\@tempb{\linewidth}{\wd\pandoc@box}%
  \ifdim\@tempb\p@<\@tempa\p@\let\@tempa\@tempb\fi% select the smaller of both
  \ifdim\@tempa\p@<\p@\scalebox{\@tempa}{\usebox\pandoc@box}%
  \else\usebox{\pandoc@box}%
  \fi%
}
% Set default figure placement to htbp
\def\fps@figure{htbp}
\makeatother
\setlength{\emergencystretch}{3em} % prevent overfull lines
\providecommand{\tightlist}{%
  \setlength{\itemsep}{0pt}\setlength{\parskip}{0pt}}
\usepackage{bookmark}
\IfFileExists{xurl.sty}{\usepackage{xurl}}{} % add URL line breaks if available
\urlstyle{same}
% fallback for those not using the hyperref driver hyperxmp:
\makeatletter
\@ifundefined{xmpquote}{\newcommand{\xmpquote}[1]{#1}}{}
\makeatother
\hypersetup{
  hidelinks,
  pdfcreator={LaTeX via pandoc}}

\author{}
\date{}

\begin{document}

{
\setcounter{tocdepth}{2}
\tableofcontents
}
\section{WebCoder --- Arhitectura
sistemului}\label{webcoder--arhitectura-sistemului}

\begin{quote}
Rezumat într-un paragraf: WebCoder este o platformă bilingvă (RO/EN +
încă 10 limbi de interfață) de programare competitivă --- un front-end
Next.js 16 peste un API Django 6 / DRF, un evaluator (judge) pe Celery
care execută cod neverificat într-un sandbox Docker întărit, printr-un
socket-proxy de filtrare, și un strat de stare Postgres+Redis. Ghidul
este pentru inginerii care trebuie să înțeleagă, să ruleze sau să
extindă sistemul; fiecare secțiune se leagă de cod pe care îl poți
deschide azi.
\end{quote}

\textbf{Cui se adresează:} contributori de backend/frontend, operatori
care deployază stack-ul și oricine evaluează designul platformei.
Perechea în engleză:
\href{../guide-en/ARCHITECTURE.md}{\texttt{docs/guide-en/ARCHITECTURE.md}}.

\begin{center}\rule{0.5\linewidth}{0.5pt}\end{center}

\subsection{1. Harta componentelor}\label{1-harta-componentelor}

\pandocbounded{\includegraphics[keepaspectratio,alt={Harta
componentelor}]{C:/Users/Prekzursil/.zcode/workspace/draft-9cac9da0-09ed-462a-8bd2-a78591461ba4/wt-polish/docs/dist/ARCHITECTURE/ro/8501c0f0ae9c3fd3ececf6d75972d736767b0077.png}}

\emph{Sursa:
\href{../diagrams/component-map.mmd}{\texttt{docs/diagrams/component-map.mmd}}}

{\def\LTcaptype{none} % do not increment counter
\begin{longtable}[]{@{}lll@{}}
\toprule\noalign{}
Componentă & Unde & Responsabilitate \\
\midrule\noalign{}
\endhead
\bottomrule\noalign{}
\endlastfoot
Aplicația Next.js 16 & rădăcina repo-ului \texttt{src/}, App Router &
pagini SSR/ISR, editoarele CodeMirror/Blockly, UI RO/EN cu i18next +
câmpuri de problemă pe 12 limbi \\
Django 6 + DRF & \texttt{backend/} (\textasciitilde25 aplicații) & REST
\texttt{/api/v1}: probleme, submissions, concursuri, ratinguri, clasă,
forum, mesagerie, traduceri, API public \\
PostgreSQL 16 & compose \texttt{postgres} & sistemul de referință (fără
port publicat pe host) \\
Redis & compose \texttt{redis} & broker celery, cache de throttle, coada
judge \\
celery-worker & compose \texttt{celery-worker} & cozile judge + default;
\texttt{acks\_late} cu garda anti-blocaj \\
judge-runner & \texttt{backend/Dockerfile.judge} & toolchain-uri
sandboxate: gcc/g++, OpenJDK, CPython, Node, fpc, mono, golang \\
socket-proxy & overlay-ul implicit compose & filtru API docker --- doar
containers/exec POST; restul 403 \\
celery-beat & compose \texttt{celery-beat} & recompute noaptea
clasamentului, purgare notificări, retenție outputuri, digest
săptămânal, \texttt{contest-tick} pe minut \\
db-backup & sidecar compose & \texttt{pg\_dump\ -Fc} zilnic, retenție
7z/4s/6l \\
\end{longtable}
}

Contractul frontend--backend: cookie-urile merg cu
\texttt{credentials:\ \textquotesingle{}include\textquotesingle{}};
lista catalogului e cache-uită ISR 60 s server-side (dezactivat sub
\texttt{E2E\_API\_BASE}).

\subsection{2. Pipeline-ul de evaluare}\label{2-pipeline-ul-de-evaluare}

\pandocbounded{\includegraphics[keepaspectratio,alt={Pipeline
evaluator}]{C:/Users/Prekzursil/.zcode/workspace/draft-9cac9da0-09ed-462a-8bd2-a78591461ba4/wt-polish/docs/dist/ARCHITECTURE/ro/43ae7a19bbb00e46ba5e5a0f957026d485dd035b.png}}

\emph{Sursa:
\href{../diagrams/judge-pipeline.mmd}{\texttt{docs/diagrams/judge-pipeline.mmd}}}

Invarianți-cheie (toți enforce-ați în cod,
\texttt{backend/submissions/}):

\begin{enumerate}
\def\labelenumi{\arabic{enumi}.}
\tightlist
\item
  \textbf{Registrul de limbaje}
  (\texttt{submissions/judge\_utils/languages.py}) e sursa unică de
  adevăr --- comenzi de compilare/rulare, multiplicatori de timp per
  limbaj (ex. python ×3), praguri de memorie, pin-uri de toolchain
  publicate pe \texttt{/api/v1/public/languages/}.
\item
  \textbf{Orice rulare judecată} e \texttt{-\/-network\ none} cu mount
  read-only al submissionului, toate capabilitățile eliminate, uid 1000.
\item
  \textbf{Scara verdicturilor} per test: TLE (wall vs CPU), MLE (scara
  max-RSS incl. exit-137 + ≥0,9×limită), RE (mesaje mapate pe semnal),
  WA, AC --- trigger-urile de fabricare a verdicturilor au fost șterse
  în valul judge-fairness.
\item
  \textbf{Moduri de comparare}: EXACT / STRIP / LINES\_STRIP / TOKEN /
  epsilon-zecimal, plus checker-e proprii cu protocolul parțial exit-7
  \texttt{POINTS\ \textless{}0..1\textgreater{}} (compatibil testlib).
\item
  \textbf{Moduri de punctaj}: BINARY, PARTIAL\_PER\_TEST, SUBTASKS
  (regula DMOJ --- toate testele grupului).
\item
  \textbf{Garanarea feedbackului}: testele demo expun input+expected
  pentru diff; testele ascunse stripă outputurile pentru privitorii
  neprivilegiați; barele de timp/memorie aplicate se persistă pe fiecare
  rând de rezultat.
\item
  \textbf{Biblioteci custom}: pip install-ul rulează într-un container
  one-shot pe rețeaua internă \texttt{judge-build} al cărei unic egress
  e un proxy allowlist PyPI pre-populat --- fail-closed implicit
  (RUNBOOK §judge-pip-egress).
\end{enumerate}

\subsection{3. Modelul de date (felia de
bază)}\label{3-modelul-de-date-felia-de-bazux103}

\pandocbounded{\includegraphics[keepaspectratio,alt={Model de
date}]{C:/Users/Prekzursil/.zcode/workspace/draft-9cac9da0-09ed-462a-8bd2-a78591461ba4/wt-polish/docs/dist/ARCHITECTURE/ro/22ac648fe9cee6cdb213cb70532bd38ea88a61b6.png}}

\emph{Sursa:
\href{../diagrams/data-model.mmd}{\texttt{docs/diagrams/data-model.mmd}}}

Note de modelare:

\begin{itemize}
\tightlist
\item
  Câmpurile localizate sunt obiecte JSON cheiate pe limbă, cu
  \textbf{engleza obligatorie} (regula owner (b) din
  \texttt{problems/i18n\_fields.py}); româna e limba pieței principale;
  valul i18n-10 a adăugat încă 10 limbi de interfață pentru
  enunțuri/titluri.
\item
  \texttt{Submission} → \texttt{SubmissionTestResult} poartă verdictul,
  timpul, memoria per test \textbf{și limitele efectiv aplicate}
  (chitanțele judge-fairness).
\item
  \texttt{Problem.reference\_solution\_code} alimentează hack-urile,
  scorul de putere al testelor și generatorul de suite ascunse
  (\texttt{boost\_hidden\_tests}).
\item
  Proveniența importurilor (\texttt{source\_url},
  \texttt{source\_authors}, ledger-ul \texttt{ImportSource})
  înregistrează de unde a venit conținutul și în ce condiții (CNT-2).
\end{itemize}

\subsection{4. Autentificarea și fluxul
cookie-urilor}\label{4-autentificarea-ux219i-fluxul-cookie-urilor}

\pandocbounded{\includegraphics[keepaspectratio,alt={Fluxul
cookie-urilor}]{C:/Users/Prekzursil/.zcode/workspace/draft-9cac9da0-09ed-462a-8bd2-a78591461ba4/wt-polish/docs/dist/ARCHITECTURE/ro/cc037c783b21e881ab154a672178aa2c479382a7.png}}

\emph{Sursa:
\href{../diagrams/auth-cookie-flow.mmd}{\texttt{docs/diagrams/auth-cookie-flow.mmd}}}

\begin{itemize}
\tightlist
\item
  \texttt{wc\_access} (5 min) + \texttt{wc\_refresh} (24 h) sunt
  \textbf{cookie-uri httpOnly} --- niciodată în \texttt{localStorage};
  rotirea blacklistează refresh-ul vechi.
\item
  Metodele unsafe cer header-ul CSRF double-submit
  (\texttt{GET\ /api/v1/auth/csrf/} eliberează cookie-ul).
\item
  Clienții non-browser folosesc Bearer + token-uri API personale (API
  public, 60/min; cu token 240/min).
\item
  Roluri: BASIC, CREATOR, TEACHER, VERIFIER, ADMIN; profilurile
  minorilor sunt ne-publice implicit (LEG-2).
\end{itemize}

\subsection{5. Ciclul de viață al
concursurilor}\label{5-ciclul-de-viaux21bux103-al-concursurilor}

\pandocbounded{\includegraphics[keepaspectratio,alt={Ciclul
concursurilor}]{C:/Users/Prekzursil/.zcode/workspace/draft-9cac9da0-09ed-462a-8bd2-a78591461ba4/wt-polish/docs/dist/ARCHITECTURE/ro/4a43b15d698febc064a1849ebdfbbe41130eb965.png}}

\emph{Sursa:
\href{../diagrams/contest-lifecycle.mmd}{\texttt{docs/diagrams/contest-lifecycle.mmd}}}

\begin{itemize}
\tightlist
\item
  \texttt{contest\_tick} (beat, pe minut) avancează tranzițiile;
  formatele sunt clase-strategie (IOI sum-of-best implicit; există ICPC
  penalty și stil AtCoder).
\item
  Freeze-ul e un filtru pe scoreboard; ratingurile se comită
  post-concurs cu Glicko + suprimarea de afișare pentru începători (stil
  AtCoder).
\item
  Upsolvingul/participarea virtuală redeschide problemele sub un ceas
  personal; hack-urile se deblochează după rezolvarea completă a tablei.
\end{itemize}

\subsection{6. Mecanici transversale}\label{6-mecanici-transversale}

\begin{itemize}
\tightlist
\item
  \textbf{Paginare}: toate endpoint-urile de listă folosesc paginare
  standard/cursor.
\item
  \textbf{Throttle}: anon 60/min, user 240/min, \texttt{auth} 10/min,
  \texttt{submit} 20/h, scope-uri separate pentru
  scratch/guest/sample-run; contoare în redis (partajate între procese).
\item
  \textbf{Observabilitate}: \texttt{/healthz}, \texttt{/readyz}
  (db+cache+broker), snapshot de metrici admin, loguri JSON
  (\texttt{LOG\_JSON=1}).
\item
  \textbf{Reproductibilitate}: \texttt{export\_catalog} /
  \texttt{import\_catalog} (+ snapshot ca asset de release)
  reconstruiesc determinist catalogul importat; \texttt{seed\_problems}
  reconstruiește pachetul de start; \texttt{boost\_hidden\_tests}
  regenerează suite ascunse din specificații comise în repo.
\item
  \textbf{Testare}: pytest (profil CI sqlite, 100 \% acoperire de linii,
  zero warninguri), vitest, Playwright mock + mod real-stack
  (\texttt{E2E\_API\_BASE}).
\end{itemize}

\subsection{7. Structura
repository-ului}\label{7-structura-repository-ului}

\begin{Shaded}
\begin{Highlighting}[]
\NormalTok{backend/\textless{}app\textgreater{}/         aplicații Django (problems, submissions, contests, ratings,}
\NormalTok{                       classroom, forum, messaging, translations, …)}
\NormalTok{backend/problems/data/ specificații de conținut comise + fixture{-}e pentru catalog}
\NormalTok{docs/                  suita de documentație (ghiduri EN/RO, dist/)}
\NormalTok{e2e/                   specificații Playwright (mock{-}api + real{-}stack)}
\NormalTok{mobile/                shell{-}ul PWA (service worker, pachete offline)}
\NormalTok{scripts/               scripturi ops (load smoke, autoscale, restore drill)}
\NormalTok{src/                   aplicația Next.js 16 (App Router)}
\NormalTok{docker{-}compose*.yml    stack{-}ul Mode{-}A + overlay{-}uri (socket{-}proxy, custom{-}libs)}
\end{Highlighting}
\end{Shaded}

\subsection{8. Lectură
suplimentară}\label{8-lecturux103-suplimentarux103}

\begin{itemize}
\tightlist
\item
  RUNBOOK (\texttt{docs/RUNBOOK.md}) --- variabile de mediu, stack-ul
  live, backupuri, modurile e2e
\item
  SECURITATE (\texttt{SECURITATE.md}) --- chitanțele sandbox-ului,
  throttle-urile, raportarea
\item
  CONTRIBUIE --- porțile de calitate, convențiile PR, fluxul valurilor
\item
  PROGRAMARE PE BLOC --- linia Blockly/Scratch și problemele cu cod
  parțial
\end{itemize}

\end{document}
